% Generated draft: requires mathematical and source review.
\documentclass[11pt]{article}
\usepackage[margin=25mm]{geometry}
\usepackage{fontspec}
\setmainfont{Latin Modern Roman}
\usepackage{amsmath,amssymb,mathtools,unicode-math}
\setmathfont{Latin Modern Math}
\usepackage{enumitem,needspace,xurl,hyperref}
\setlength{\parindent}{0pt}
\setlength{\parskip}{5pt}
\setlength{\emergencystretch}{3em}
\newcommand{\N}{\mathbb{N}}
\newcommand{\Z}{\mathbb{Z}}
\newcommand{\Q}{\mathbb{Q}}
\newcommand{\R}{\mathbb{R}}
\newcommand{\C}{\mathbb{C}}
\newcommand{\F}{\mathbb{F}}
\newcommand{\A}{\mathbb{A}}
\newcommand{\PP}{\mathbb{P}}
\newcommand{\E}{\mathbb{E}}
\newcommand{\eps}{\varepsilon}
\newcommand{\dd}{\,\mathrm{d}}
\begin{document}
\section*{UnsolvedMath ID: 30005970}
\textbf{Source problem code:} catalogue entry
\subsection{Definitions and mathematical statement}
Let \(N\) be a positive integer, and let \(X_0(N)\) be the compactified modular curve over \(\Q\) parametrizing elliptic curves with a cyclic subgroup of order \(N\). Let \(j:X_0(N)\to X(1)\cong\PP^1_{\Q}\) be the usual morphism forgetting that subgroup. The source requires \(N\) to be ``large'', but supplies neither a bound nor any further restrictions on \(N\). Write
\[
G_N=\mathrm{PGL}_2(\Z/N\Z)
=\mathrm{GL}_2(\Z/N\Z)\big/\{uI:u\in(\Z/N\Z)^\times\}.
\]
For a closed point \(P\) of a curve over \(\Q\), its degree is \([\Q(P):\Q]\). A conventional interpretation of its Galois group is the Galois group of the normal closure of \(\Q(P)/\Q\), acting on the embeddings of \(\Q(P)\) into \(\overline{\Q}\). Such an action is primitive if it is transitive and preserves no nontrivial partition into blocks. The source does not explicitly specify whether its comparisons of Galois groups are abstract group comparisons or comparisons of permutation groups.

Let \(U_N\subseteq X(1)\) be the open subset over which \(j\) is finite and étale. For \(\alpha\in U_N(\Q)\), let \(\Gamma_{N,\alpha}\) be the image of the absolute Galois group of \(\Q\) acting on the geometric points of \(j^{-1}(\alpha)\); equivalently, it is the Galois group of the splitting field of that finite étale fiber, with its fiber action. This definition makes sense even when the fiber is not a single closed point. A thin subset of \(X(1)(\Q)\) is a subset contained in a finite union of a proper closed subset's rational points and images of rational points under dominant generically finite maps of degree at least two from integral varieties of the same dimension.

The supplied formulation asserts that Hilbert irreducibility provides a thin set \(S\subseteq X(1)(\Q)\) outside which the fiber has Galois group \(G_N\). This assertion is retained as supplied context: the necessary restrictions on \(N\), the generic Galois group, and the intended permutation action are not established by the supplied excerpt. Its questions are whether \(S\) is finite or infinite and, if infinite, which Galois groups occur infinitely often among exceptional fibers. However, an exceptional thin set witnessing Hilbert irreducibility is not unique: it can be enlarged by another thin set. Consequently, these questions require a specified choice of \(S\).

A precise candidate for the intended exceptional set, using abstract group isomorphism and excluding branch values, is
\[
E_N=\{\alpha\in U_N(\Q):\Gamma_{N,\alpha}\not\cong G_N\}.
\]
If \(S\) is intended to mean this set, the questions become: for each admissible fixed \(N\), is \(E_N\) finite; and, when it is infinite, for which finite groups \(H\not\cong G_N\) is
\[
\{\alpha\in E_N:\Gamma_{N,\alpha}\cong H\}
\]
infinite? This is an explicitly qualified reformulation, since the source does not identify \(S\) with \(E_N\), specify the admissible levels, or explain the treatment of branch values.

The source also asks whether there is an infinite family of algebraic points on \(X_0(N)\) with Galois group \(G_N\) that ``does not arise from \(j\)''. This last condition is undefined in the record. If it means that the points lie outside fibers above rational values, it would impose \(j(P)\notin X(1)(\Q)\); this is only a possible interpretation. The intended degrees, permutation representations, and meaning of a family also remain unspecified.

Finally, the broader requested extension concerns a result attributed to Khawaja--Siksek, labeled ``Khawaja--S.'' in the excerpt: under omitted ``same hypotheses'', infinitely many degree-\(n\) points with Galois group \(S_n\) or \(A_n\) imply that only finitely many degree-\(n\) points have a primitive Galois group different from both. The problem asks for extensions to other families of simple and almost simple groups. A simple group here is a nontrivial group with no proper nontrivial normal subgroup; a finite group \(G\) is almost simple if \(T\leq G\leq\operatorname{Aut}(T)\) for some nonabelian finite simple group \(T\). The missing theorem hypotheses and unspecified replacement group families prevent a complete quantified extension statement.
\subsection{Short English statement}
The problem asks whether exceptional rational fibers of the modular \(j\)-map are finite, which exceptional Galois groups recur infinitely often, and whether infinitely many points with the stated full Galois group arise in another way. It also seeks extensions of a Khawaja--Siksek finiteness theorem. The level restrictions, the exact exceptional set, the meaning of ``arise from \(j\)'', and the hypotheses of the cited theorem require clarification. The supplied open status is dated curation, not an independently verified conclusion.
\subsection{Sources}
\begin{itemize}
\item Local UnsolvedMath registry, supplied complete record for permanent ID \(30005970\), titled ``Exceptional Galois Points on Modular Curves'', exact source problem code \texttt{catalogue entry}, in the Oberwolfach Reports Open Problems collection. The registry supplies proposed year \(2024\) and names no proposer.
\item Supplied report citation: \emph{Explicit Methods in Number Theory (2025)}, Oberwolfach Reports, pp.~2309--2368, DOI \texttt{10.4171/owr/2024/40}. Supplied source URL and literature-source link labeled ``original OWR report'': \url{https://doi.org/10.4171/owr/2024/40}. No independent consultation of this report is claimed. The citation's \(2025\) designation and the DOI's \(2024\) component are retained as supplied.
\item The supplied excerpt attributes Theorem~3 to ``Khawaja--S.''; the cleaned registry statement identifies the attribution as Khawaja--Siksek. Its reference to ``the same hypotheses'' has no antecedent in the supplied excerpt. No separate complete citation for that theorem is supplied.
\item Literature assessment checked at exactly \texttt{2026-08-22}: the supplied curation labels the problem open and describes finiteness of the thin exceptional set and existence of a family not arising from \(j\) as open questions in the primary report. This is supplied dated curation, not independently verified literature status. The background explicitly requests independent verification of time-sensitive status and cited claims.
\item The registry labels the statement \texttt{corrected\_verified}, explaining that only formula delimiters, notation, grammar, and typography were repaired while mathematical content and scope were preserved. This label does not establish mathematical correctness or resolve the missing qualifications.
\end{itemize}
\end{document}
